%*************************************
% บทที่ 1
%*************************************
\chapterTitle{1}{บทนำ}

%*******************************************
% ความเป็นมาและความสำคัญ
%*******************************************
\section{ความเป็นมาและความสำคัญ} 

%ระบุความเป็นมาและความสำคัญของโปรเจ็ค
\hspace*{1.3em} %ย่อหน้า
การเรียนการสอนในภาควิชาเทคโนโลยีสารสนเทศมีความจำเป็นต้องใช้ระบบปฏิบัติการลีนุกซ์ในหลายรายวิชา โดยเฉพาะในรายวิชาที่เกี่ยวข้องกับระบบเซิร์ฟเวอร์ การพัฒนาแอปพลิเคชัน และการจัดการโครงสร้างพื้นฐานด้านเทคโนโลยีสารสนเทศ อย่างไรก็ตาม การเตรียมสภาพแวดล้อมสำหรับการใช้งานระบบลีนุกซ์นั้นมักมีความยุ่งยากและใช้เวลานาน ทั้งในด้านการติดตั้ง การตั้งค่า และการจัดการทรัพยากร ส่งผลให้การเรียนการสอนไม่สามารถดำเนินไปได้อย่างเต็มประสิทธิภาพ

\hspace*{1.3em}
ที่ผ่านมา นักศึกษาและอาจารย์ต้องใช้วิธีการจำลองระบบลีนุกซ์ผ่านเครื่องเสมือน (Virtual Machine), ระบบย่อยของวินโดว์สำหรับลีนุกซ์ (WSL), หรือบริการคลาวด์จากภายนอก ซึ่งแต่ละวิธีมีข้อจำกัด เช่น ความยุ่งยากในการตั้งค่า ความไม่เสถียรของระบบ หรือค่าใช้จ่ายเพิ่มเติมในการใช้งานบริการจากผู้ให้บริการภายนอก โดยเฉพาะในการทำโครงงานที่ต้องใช้เซิร์ฟเวอร์ นักศึกษามักต้องจัดหาเครื่องคอมพิวเตอร์ส่วนตัวหรือเซิร์ฟเวอร์ด้วยตนเอง ซึ่งเป็นภาระทั้งด้านเวลาและงบประมาณ

\hspace*{1.3em}
เพื่อแก้ไขปัญหาดังกล่าวข้างต้น โครงงานนี้จึงมีแนวคิดในการพัฒนาแพลตฟอร์มคลัสเตอร์ภายในภาควิชา โดยใช้เทคโนโลยี Virtualization ร่วมกับระบบ cloud-Init เพื่อสร้างเครื่องจำลองลีนุกซ์ที่สามารถใช้งานได้อย่างรวดเร็วและง่ายดาย พร้อมทั้งมีระบบ Web Application สำหรับการจัดการและเข้าใช้งานเครื่องเสมือน ทั้งในส่วนของนักศึกษาและอาจารย์



%*******************************************
% วัตถุประสงค์ของโครงงาน
%*******************************************
\section{วัตถุประสงค์ของโครงงาน}

%ระบุวัตถุประสงค์ของโปรเจ็ค
\begin{mycustomenum}[label=1.2.\arabic*] %เพิ่มวัตถุประสงค์หลัง \item ถ้ามีหลายข้อ ก็เพิ่มบรรทัด \item ไปได้เรื่อย ๆ
    \item เพื่อออกแบบและพัฒนาแพลตฟอร์มคลัสเตอร์เสมือน (Virtual Cluster Platform) ที่สามารถสร้างและจัดการเครื่องเสมือนลีนุกซ์ได้อย่างรวดเร็วและมีประสิทธิภาพ สำหรับสนับสนุนการเรียนการสอนและการทำโครงงานในภาควิชาเทคโนโลยีสารสนเทศ
    \item เพื่อพัฒนา Web Application แบบ Role-Based Access Control (RBAC) ที่รองรับการใช้งานของทั้งนักศึกษาและอาจารย์ ในการร้องขอ จัดการ และเข้าถึงเครื่องเสมือน โดยคำนึงถึงความง่ายในการใช้งาน เสถียรภาพของระบบ และความปลอดภัยในการเข้าถึงแพลตฟอร์ม
    \newpage
    \item เพื่อเพิ่มประสิทธิภาพการใช้ทรัพยากรเซิร์ฟเวอร์ที่มีอยู่ภายในภาควิชา ผ่านการพัฒนาระบบคลัสเตอร์แบบรวมศูนย์ ที่สามารถติดตาม วิเคราะห์ และบริหารจัดการทรัพยากรได้อย่างเหมาะสม โดยมีระบบ Monitoring และ Cloud Storage ที่สามารถรองรับการใช้งานในลักษณะ multi-user environment
\end{mycustomenum}



%*******************************************
% ขอบเขตของการทําโครงงานพิเศษ
%*******************************************
\section{ขอบเขตของการทําโครงงาน} \label{sec1:scope}
\hspace*{1 cm}\textbf{ภาคการศึกษาที่ 1/2568}
%ระบุขอบเขตของโปรเจ็ค
\begin{mycustomenum2}
    \item จัดทำระบบการจัดการเครื่องคอมพิวเตอร์เซิร์ฟเวอร์
    \begin{mycustomenum2}
        \item จัดทำระบบ Hypervisor สำหรับการทำ Virtualization
        \item ระบบการจัดการเครื่องเสมือนที่อยู่ภายใต้ระบบ Hypervisor ด้วย API
        \begin{mycustomenum2}
            \item การสร้างเครื่องเสมือนบนระบบ Hypervisor
            \item การแก้ไขข้อมูลเครื่องเสมือนบนระบบ Hypervisor
            \item การลบเครื่องเสมือนบนระบบ Hypervisor
            \item การจัดการสถานะเครื่องเสมือนบนระบบ Hypervisor
        \end{mycustomenum2}
    \end{mycustomenum2}
    
    \item จัดทำระบบ Web Application สำหรับการเข้าใช้งานระบบ
    \begin{mycustomenum2}
        \item ระบบเข้าใช้งานสำหรับนักศึกษา
        \begin{mycustomenum2}
            \item ระบบการยืนยันตัวตนด้วย Google Account เพื่อเข้าใช้งานระบบด้วยโดเมน @email.kmutnb.ac.th และยืนยันภาควิชาที่นักศึกษาสังกัดด้วยรหัสนักศึกษาที่ปรากฎบนที่อยู่อีเมลล์
            \item ระบบยื่นคำร้องขอการใช้งานเครื่องเสมือนสำหรับการใช้งานในการเรียนปกติหรือสำหรับการทำโครงงานพิเศษ
            \item ระบบการจัดการเครื่องเสมือนและการเข้าใช้งานเครื่องเสมือนของนักศึกษาที่ยังมีสถานะปกติ
            \item ระบบยื่นคำร้องขอต่ออายุการใช้งานเครื่องเสมือนที่มีอยู่เพื่อขอใช้งานเครื่องเสมือนต่อในภาคการศึกษาถัดไป
        \end{mycustomenum2}            
        \item ระบบเข้าใช้งานสำหรับอาจารย์หรือผู้ดูแลระบบ
        \begin{mycustomenum2}
            \item ระบบการยืนยันตัวตนด้วย Google Account เพื่อเข้าใช้งานระบบด้วยโดเมน @itm.kmutnb.ac.th และจำกัดการเข้าใช้ระบบเฉพาะบุคลากรภายในภาควิชา
            \item ระบบจัดการคำร้องขอใช้งานเครื่องเสมือนและคำร้องขอต่ออายุการใช้งานเครื่องเสมือนของนักศึกษา
            \item ระบบการสร้างเครื่องเสมือนสำหรับการใช้งานของอาจารย์หรือผู้ดูแลระบบ
            \item ระบบการจัดการและเข้าใช้งานเครื่องเสมือนของอาจารย์หรือผู้ดูแลระบบ
            \item ระบบการจัดการเครื่องเสมือนของนักศึกษาในระบบ โดยสามารถเปลี่ยนแปลงสถานะของเครื่องเสมือนให้อยู่ในรูปแบบการจัดเก็บเครื่องเสมือนถาวรหรือการลบเครื่องเสมือนของนักศึกษา
            \item ระบบการกำหนดช่วงเวลาของภาคการศึกษา เพื่อให้ระบบสามารถจัดหมวดหมู่ของเครื่องเสมือนในภาคการศึกษานั้น ๆ ได้
            \item ระบบการจัดการรายชื่ออาจารย์หรือผู้ดูแลที่สังกัดภาควิชา
        \end{mycustomenum2}  
    \end{mycustomenum2}
\end{mycustomenum2}

\hspace*{1 cm}\textbf{ภาคการศึกษาที่ 2/2568}
%ระบุขอบเขตของโปรเจ็ค
\begin{mycustomenum2}
    \item จัดทำระบบ Reverse Proxy สำหรับการเข้าใช้งานพอร์ตต่าง ๆ ของเครื่องเสมือนจากระบบเครือข่ายภายสาธารณะได้
    \item จัดทำระบบ Cloud Storage สำหรับการจัดเก็บไฟล์สำหรับการใช้งานของอาจารย์หรือผู้ดูแลระบบภายในภาควิชา
    \item จัดทำระบบ Monitoring ที่จัดเก็บข้อมูล Trace, Metric และ Log โดยสามารถแสดงผลผ่าน Dashboard ได้
    \item จัดทำการระบบเพิ่มเติมบน Web Application
    \begin{mycustomenum2}
        \item ระบบการกำหนดค่า Port Forwarding บนระบบ Reverse Proxy เพื่อกำหนดพอร์ตที่ต้องการใช้งานจากระบบเครือข่ายภายนอก
        \item ระบบการเข้าใช้งาน Cloud Storage สำหรับอาจารย์หรือผู้ดูแลระบบ เพื่อใช้งานในการจัดเก็บไฟล์
    \end{mycustomenum2}
    \item ประเมินประสิทธิภาพของระบบ
    \begin{mycustomenum2}
        \item วัดค่า Provisioning Time วัดค่าเฉลี่ย (Mean) และ ค่ามัธยฐาน (Median) ของเวลาที่ใช้ในการสร้างเครื่องเสมือนใหม่ ตั้งแต่เริ่มต้นกระบวนการจนพร้อมใช้งานเปรียบเทียบกับวิธีการเดิม
        \item วัดค่า Concurrent User Performance ทดสอบความเร็วการตอบสนองของระบบเมื่อมีผู้ใช้งานพร้อมกันหลายราย วัดค่าการตอบสนองของระบบ (Response Time) เฉลี่ย (ms) และอัตราความล้มเหลว (Error Rate) (\%)
    \end{mycustomenum2}
\end{mycustomenum2}



%*******************************************
% ประโยชน์ที่คาดว่าจะได้รับ
%*******************************************
\section{ประโยชน์ที่คาดว่าจะได้รับ}

%ระบุประโยชน์ที่คาดว่าจะได้รับของโปรเจ็ค
\begin{mycustomenum}[label=1.4.\arabic*] %เพิ่มวัตถุประสงค์หลัง \item ถ้ามีหลายข้อ ก็เพิ่มบรรทัด \item ไปได้เรื่อย ๆ
    \item ได้ระบบแพลตฟอร์มคลัสเตอร์ที่สามารถใช้งานได้จริงภายในภาควิชาเทคโนโลยีสารสนเทศ เพื่อสนับสนุนการเรียนการสอนและการทำโครงงานของนักศึกษาและอาจารย์
    \item ได้ระบบเครื่องเสมือนลีนุกซ์ที่สามารถสร้างและจัดการได้อย่างรวดเร็วผ่าน Web Application โดยไม่ต้องติดตั้งระบบปฏิบัติการเพิ่มเติมบนเครื่องของผู้ใช้งาน
    \newpage
    \item ได้ใช้ประโยชน์จากเครื่องเซิร์ฟเวอร์ที่มีอยู่ในภาควิชาอย่างเต็มประสิทธิภาพ โดยนำมาให้บริการในรูปแบบคลัสเตอร์สำหรับการเรียนการสอน
    \item ได้โครงงานที่สามารถนำไปใช้งานจริงภายในภาควิชา และสามารถต่อยอดพัฒนาเพิ่มเติมในอนาคต
    \item ได้โครงงานอื่น ๆ ที่สามารถใช้งานได้จริงในภาควิชา โดยนำระบบที่นักศึกษาที่พัฒนาในโครงงานมาติดตั้งบนแพลตฟอร์มคลัสเตอร์
\end{mycustomenum}
